\documentclass[11pt]{article}
\usepackage{amsmath,amssymb,amsthm}
\usepackage[margin=1in]{geometry}
\newtheorem{theorem}{Theorem}
\newtheorem{lemma}[theorem]{Lemma}
\newtheorem{proposition}[theorem]{Proposition}
\newtheorem{conjecture}[theorem]{Conjecture}
\theoremstyle{remark}
\newtheorem{remark}[theorem]{Remark}
\newcommand{\Q}{\mathbb Q}
\newcommand{\Z}{\mathbb Z}
\newcommand{\C}{\mathbb C}
\title{P4, Room of Mathematicians, Seat: Erd\H{o}s --- a search for a
non-pole-counting obstruction to $D$-algebraicity of $U$}
\date{2026-09-26}
\begin{document}
\maketitle

\noindent Labels as in \texttt{P4\_note.tex}: PROVED / VERIFIED / HEURISTIC / OPEN.

\section*{Calibration point}

I have read \texttt{code/zeta\_probe/rootunity\_zeros/general/P4/P4\_note.tex} in full, reviewed
version, correction included. My understanding, restated to confirm it: the question is whether
$U(x)=\sum u_nx^n$ (paper2a's growth series, $=\mathcal W(x,x^2)$ in the model (M)) satisfies
\emph{some} algebraic differential equation $P(x,U,U',\ldots,U^{(k)})=0$, $P\not\equiv0$ a
polynomial. The note establishes: (i) density of poles on a positive-measure boundary arc is
\emph{not} by itself an obstruction, because nonlinear ODEs can have movable Painlev\'e-type poles
dense on a curve (Remark~\texttt{rem:nogeneralize}); the linear-ODE fixed-singularity argument
(Lemma~\texttt{lem:odeposes}) has no literal nonlinear analogue; (ii) the attempted elliptic-pullback
counterexample was arithmetically broken and is retracted; (iii) the correct tool family is
Nevanlinna-theoretic (Malmquist/Steinmetz/Bank--Kaufman), which needs a disc-analogue growth theorem
that is not confirmed to exist in citable form, plus a quantitative lower bound on the pole-counting
function $n(r,U)$, neither of which is established. My assignment is to look for a
\emph{completely different} elementary obstruction, not resting on pole-counting or growth theorems
at all, and explicitly not to touch pole-counting (Tao's seat) or redo the elliptic construction
(already correctly retracted).

\section{What I tried}

\subsection{Idea 1: functional-equation degree incompatibility}

$U=\mathcal W(x,x^2)$ where $\mathcal W(x,y)$ is built (Theorem~\texttt{thm:model}, (M3)) from a
resolvent $(I-\mathcal T(q))^{-1}$ of a compact linear operator $\mathcal T(q)$ on $\ell^1$, together
with a scalar self-referential quantity $B_0=1/(1-\mathfrak g(q,y)t_1(q))$, $t_1=v^\top(I-M_0)^{-1}u$,
where $M_0$ is another compact operator built the same way. This is genuinely a \emph{linear}
structure (Neumann series of compact operators), not a nonlinear functional/Mahler-type equation in
$x$ of the kind that has a useful ``degree blows up under iteration'' argument (that trick applies to
equations like $U(x)=R(x,U(qx))$ with $R$ a fixed rational map, where substituting $U(qx)=R(qx,U(q^2x))$
etc.\ and comparing against a hypothetical fixed-order/degree ODE forces the degree to grow with the
number of substitutions, a contradiction). Paper2a's own strongest results in exactly this style
(\texttt{prop:noqdiff}: no linear $q$-difference or Mahler equation; \texttt{cor:Qminusone}: no
$Q=-1$ relation) already show $U$ satisfies \emph{no} equation of the substitutable form $R(x,U(qx))$
at all --- linear or otherwise, that is presumably the content that would be needed to feed a
degree-counting argument. So there is no functional equation of the shape idea 1 needs to attack:
the input a Mahler/functional-degree argument requires is exactly the thing paper2a already rules out
existing, which blocks the argument rather than feeding it. I could not construct a substitute
functional relation for $U$ itself (as opposed to for the operator-valued objects $\mathcal T(q)$,
$M_0(q)$, which do satisfy the trivial linear identity $\mathcal T(q)=q\cdot(\text{fixed matrix in
}q^{\max(s,s')}\text{ form})$, but this is a statement about matrix \emph{entries}, not about $U(x)$
as a scalar function, and does not obviously transport to a constraint on $U$ under the assumption
$P(x,U,\ldots,U^{(k)})=0$). \textbf{Verdict on idea 1: does not go through}, and for a structural
reason (paper2a already proved the negative form of the needed hypothesis), not merely because I
ran out of time.

\subsection{Idea 2: algebraic independence of $U,xU',x^2U'',\ldots$}

A direct algebraic-independence argument (e.g.\ producing an explicit nonvanishing generalized
Wronskian of $U$ and finitely many of its derivatives, or a differential-Galois-theoretic
transcendence-degree count) would need either (a) an explicit closed form for $U^{(j)}$ in terms of
resolvent derivatives $\partial_q^j\bigl[(I-\mathcal T)^{-1}\bigr]$, which are themselves given only
implicitly through a Neumann series with no closed algebraic form (the operator $\mathcal T(q)$ has
infinitely many nonzero entries as $q$ varies; its resolvent is not a rational or algebraic function
of $q$ except possibly at isolated special points), or (b) a purely numerical/algebraic-independence
argument (e.g.\ a Rust-computed Taylor expansion feeding a PSLQ-type search for a polynomial relation
among a finite window of coefficients, as was done for the bedrock coefficients
[\texttt{bedrock-coeffs-irreducible.md}]). Route (b) would only ever produce evidence \emph{against}
an ODE of bounded, tested order/degree (a relation search that comes back empty for orders up to some
bound $k_{\max}$), never a proof that \emph{no} order works --- which is exactly the same
``uniform-in-$(k,d)$'' obstacle that \texttt{P4\_note.tex}\S\texttt{sec:gap}(c) identifies for the
growth-theorem route, arrived at from a different direction. I did not run this numerically: it would
not close the row even if it came back clean, and a positive hit (a spurious low-order relation found
by an under-powered search) would be actively misleading, so per the ``no forced results'' standing
instruction I did not spend the compute budget on a search whose informative value is capped from the
start. \textbf{Verdict on idea 2: same uniformity obstacle as the growth-theorem route, reached
independently; no genuinely new elementary route found.}

\subsection{Idea 3: a dimension/parameter-counting argument from Taylor-coefficient arithmetic}

If $U$ solved $P(x,U,\ldots,U^{(k)})=0$ of order $k$ with a leading coefficient not vanishing at
$x=0$, then by the standard existence/uniqueness theory for algebraic ODEs (formal Cauchy--Kovalevskaya:
differentiate $P=0$ repeatedly and solve for $U^{(k+1)},U^{(k+2)},\ldots$ algebraically in terms of
$x,U,\ldots,U^{(k)}$, valid as long as $\partial P/\partial U^{(k)}\ne0$ along the solution), the
\emph{entire} Taylor sequence $(u_n)$ is a fixed rational function of the first $k+1$ coefficients
$u_0,\ldots,u_k$ together with $x,\ldots$ evaluated recursively --- i.e.\ $(u_n)_{n>k}$ is determined,
by an explicit (if complicated) recursion of bounded ``local complexity'' per step, from finitely much
initial data. This is the precise elementary fact I could try to contradict using known arithmetic
properties of $u_n$ (they are honest combinatorial counts, non-negative integers, from paper2a's
enumeration). The trouble: ``bounded local complexity per step'' does not bound the \emph{total}
complexity of computing $u_n$ from $u_0,\ldots,u_k$ (differentiating a fixed polynomial relation $n-k$
times produces expressions of size exponential, or worse, in $n$ in general, so this gives essentially
no constraint on the growth rate or the arithmetic nature of $u_n$ without already having a
Nevanlinna-type growth bound in hand for the exponents/degrees appearing --- which is again exactly
the missing ingredient of \texttt{P4\_note.tex}). I could not find a way to convert the qualitative
fact ``$(u_n)$ is a rational function of finitely many initial values'' into a numerically checkable
contradiction using only elementary bounds (e.g.\ integrality, sign, or a simple recursion-order
argument as in \texttt{cor:Qminusone}'s $2$-adic valuation trick): the $2$-adic argument that closes
\texttt{cor:Qminusone} is special to the linear $Q=-1$ functional relation and does not have an
evident nonlinear-ODE analogue, for the same reason idea 1 fails --- there is no functional relation
of the needed shape to feed it. \textbf{Verdict on idea 3: the parameter-counting fact is true but
content-free without a growth estimate, i.e.\ it collapses onto the same missing ingredient as the
already-identified route, not a new one.}

\section{Verdict}

\textbf{NOTHING WORKED} as a genuinely independent elementary obstruction. All three candidate
directions (functional-equation degree counting, algebraic independence via Wronskians/PSLQ,
parameter counting from the Cauchy--Kovalevskaya recursion) were checked honestly and each either:
(a) needs exactly the ingredient \texttt{P4\_note.tex} already identified as missing (a quantitative,
uniform-in-order/degree growth or counting statement), reached from a different door but the same
room, or (b) is blocked structurally by results paper2a already has in the opposite direction
(\texttt{prop:noqdiff}, \texttt{cor:Qminusone} rule out exactly the kind of functional relation a
degree-counting argument would need as its hypothesis, rather than supplying one). I did not find,
and want to be explicit that I did not force, any construction that sidesteps Nevanlinna-type growth
theory entirely. This is consistent with the general shape of the D-algebraicity literature: I am not
aware of an elementary (non-growth-theoretic) technique for proving a specific, concretely given
transcendental function is \emph{not} D-algebraic, only techniques based on (i) fixed/movable
singularity structure (the linear case, already used and shown not to generalize), (ii) Nevanlinna
growth bounds (the route already identified and still open), or (iii) exhibiting an actual algebraic
relation among $U$ and its jets numerically and then trying to prove it holds (which would prove $U$
\emph{is} D-algebraic, the opposite direction, and which idea 2's numerics would only ever suggest,
never rule out, for all orders at once).

No computation was run for this note (per \S\ref{sec:idea2} above, a bounded-order relation search
would not have been decisive either way and risked a misleading spurious hit). No file besides this
one was modified. This seat's contribution to the room is negative: it closes off idea 1--3 as stated,
and confirms, from an independent elementary angle, that \texttt{P4\_note.tex}'s conclusion --- the
row stays OPEN, and the one substantive missing ingredient is a quantitative growth/counting theorem
--- is not an artifact of having looked only at pole-counting; the elementary alternatives collapse
onto the same gap.

\end{document}
